\providecommand{\tightlist}{\setlength{\itemsep}{0pt}\setlength{\parskip}{0pt}}
\newpage
\section{Algoritmos quânticos}
\subsection{Oráculos e phase kickback}
Você tem uma função \(f\) guardada numa ``caixa-preta'': dá para perguntar o valor de \(f(x)\), mas não dá para ler o código por dentro. Quantas perguntas são necessárias para descobrir uma propriedade de \(f\)? Os primeiros algoritmos quânticos respondem a esse tipo de pergunta, e o objeto que torna isso possível é o \textbf{oráculo}.

\subsubsection{Como é um algoritmo quântico}\label{como-uxe9-um-algoritmo-quuxe2ntico}

Quase todo algoritmo quântico segue o mesmo roteiro:

\begin{enumerate}
\def\labelenumi{\arabic{enumi}.}
\tightlist
\item
  \textbf{Preparar} um estado inicial simples, em geral \(\ket{0\cdots0}\).
\item
  \textbf{Aplicar portas} (um circuito) que espalham a informação em superposição e fazem as amplitudes interferirem.
\item
  \textbf{Medir} e ler um resultado clássico.
\end{enumerate}

O truque é que o passo 2 faz as amplitudes ``erradas'' se cancelarem e as ``certas'' se somarem, de modo que a medida final revela algo útil com alta probabilidade. Quando o algoritmo trabalha com uma função \(f\) dada como caixa-preta, o passo 2 inclui chamadas ao oráculo.

\subsubsection{O oráculo de bit}\label{o-oruxe1culo-de-bit}

Seja \(f:\{0,1\}^n \to \{0,1\}\) uma função que recebe \(n\) bits e devolve um bit. Uma porta quântica precisa ser reversível (unitária), e \(\ket{x}\mapsto\ket{f(x)}\) geralmente não é: se \(f\) vale \(0\) em duas entradas, não dá para voltar atrás. A saída é \textbf{guardar a entrada} e somar o resultado num qubit extra:

\[
U_f\ket{x,y} = \ket{x,\; y\oplus f(x)},
\]

onde \(x\) tem \(n\) bits, \(y\) é um bit e \(\oplus\) é a soma módulo 2 (o XOR). Usar \(y=0\) dá \(U_f\ket{x,0}=\ket{x,f(x)}\), o valor de \(f\).

\textbf{Exemplo com \(n=1\).} Seja \(f(0)=0\) e \(f(1)=1\) (ou seja, \(f(x)=x\)). Então \(U_f\ket{0,y}=\ket{0,y}\) e \(U_f\ket{1,y}=\ket{1,1\oplus y}\): o qubit \(y\) só é invertido quando \(x=1\). Isso é exatamente a porta \textbf{CNOT}. Já com \(f(x)=0\) para todo \(x\), \(U_f\) é a identidade, e com \(f(x)=1\) para todo \(x\), \(U_f=I\otimes X\).

Seja \(f(0)=1\) e \(f(1)=0\). Qual é o valor do segundo bit depois de aplicar \(U_f\) a \(\ket{0,1}\)?

\(U_f\) é \textbf{unitária} (é uma permutação dos estados da base) e é sua própria inversa, pois aplicá-la duas vezes soma \(f(x)\) duas vezes e \(f(x)\oplus f(x)=0\).

Monte a matriz de \(U_f\) para uma função dada pela tabela de valores.

\subsubsection{Phase kickback: o resultado vai para a fase}\label{phase-kickback-o-resultado-vai-para-a-fase}

Agora um truque central. Prepare o qubit extra em \(\ket{-}=\tfrac{1}{\sqrt2}(\ket0-\ket1)\) e aplique \(U_f\):

\[
U_f\ket{x}\ket{-}=\tfrac{1}{\sqrt2}\big(\ket{x,f(x)}-\ket{x,1\oplus f(x)}\big).
\]

Veja os dois casos:

\begin{itemize}
\tightlist
\item
  Se \(f(x)=0\): sobra \(\tfrac{1}{\sqrt2}(\ket{x,0}-\ket{x,1})=\ket{x}\ket{-}\), nada muda.
\item
  Se \(f(x)=1\): sobra \(\tfrac{1}{\sqrt2}(\ket{x,1}-\ket{x,0})=-\ket{x}\ket{-}\), um sinal de menos aparece.
\end{itemize}

Os dois casos cabem numa fórmula só:

\[
U_f\ket{x}\ket{-}=(-1)^{f(x)}\ket{x}\ket{-}.
\]

O qubit extra continua em \(\ket{-}\), intacto, e o valor \(f(x)\) ``voltou'' (\emph{kickback}, em inglês: coice) para o primeiro registrador como uma \textbf{fase} \(\pm1\). Como o qubit \(\ket{-}\) não muda, podemos esquecê-lo e dizer que temos um \textbf{oráculo de fase}:

\[
O_f\ket{x}=(-1)^{f(x)}\ket{x},
\]

uma matriz diagonal com \((-1)^{f(x)}\) na diagonal. Por linearidade, numa superposição cada componente \(x\) ganha o seu próprio sinal:

\[
O_f\sum_x a_x\ket{x}=\sum_x (-1)^{f(x)}a_x\ket{x}.
\]

Se \(f(x)=1\) para um certo \(x\), qual é o fator que multiplica \(\ket{x}\) no oráculo de fase?

O mesmo vale para qualquer autovetor: se \(U\ket u=e^{i\varphi}\ket u\) e o controle está em superposição, \(\ket u\) passa intacto e a fase \(e^{i\varphi}\) vai para o controle. Essa ideia reaparece na estimação de fase.

A fase global \(-1\) sozinha é invisível, mas aqui ela é \textbf{relativa}: só alguns \(x\) trocam de sinal, e é isso que a interferência vai explorar nos próximos algoritmos.

\subsubsection{Registradores e ordem dos bits (revisão)}\label{registradores-e-ordem-dos-bits-revisuxe3o}

Você já viu que \(n\) qubits guardam um número natural em binário, por exemplo \(5=0101_2\Rightarrow\ket{0101}\), e que \(n\) qubits cobrem \(2^n\) valores. Aqui vale uma \textbf{convenção}: o site escreve o bit mais significativo primeiro (\emph{big-endian}), então em \(\ket{x_1x_2\cdots x_n}\) o qubit \(x_1\) é o de maior peso e o índice na base é \(2^{n-1}x_1+\cdots+x_n\). Por exemplo, em \(\ket{x,y}\) o índice é \(2x+y\). O Qiskit usa a ordem contrária (\emph{little-endian}); cuidado ao comparar com ele.

Quantos qubits são necessários para representar todos os naturais de \(0\) a \(999\)?

\subsubsection{\texorpdfstring{Como se constrói \(U_f\)}{Como se constrói U\_f}}\label{como-se-construxf3i-u_f}

Até aqui \(f\) foi uma caixa-preta. Se tivermos a tabela de \(f\), podemos montar \(U_f\) com uma porta \(X\) multicontrolada no qubit extra para cada entrada \(x\) com \(f(x)=1\) (controles ajustados para casar com \(x\)): funciona, mas o circuito tem tamanho exponencial em \(n\). Se \(f\) tiver uma fórmula, dá para calculá-la com Toffoli e CNOT usando qubits auxiliares, que depois são desfeitos (descomputação) para não deixar lixo; aí o tamanho é linear. Veja a aula de registradores e o somador quântico. Para uma \(f\) com saída de \(m\) bits, aplicamos o mesmo raciocínio bit a bit.

\subsubsection{Complexidade de consultas e de circuito}\label{complexidade-de-consultas-e-de-circuito}

Há duas maneiras de medir o custo de um algoritmo com oráculo:

\begin{itemize}
\tightlist
\item
  \textbf{Complexidade de consultas}: quantas vezes chamamos \(U_f\) (ou \(O_f\)). Cada chamada conta como um passo, não importa o que há dentro.
\item
  \textbf{Complexidade de circuito}: quantas portas elementares o circuito usa no total, incluindo o interior do oráculo e as portas entre as chamadas.
\end{itemize}

Os algoritmos das próximas aulas mostram ganhos na complexidade de \textbf{consultas}: um algoritmo quântico pode precisar de bem menos perguntas à caixa-preta que qualquer algoritmo clássico. Isso é uma prova limpa de vantagem, mas é mais modesto do que parece: se implementarmos \(f\) de verdade, o custo do oráculo entra na conta.

\subsubsection{Resumo}\label{resumo}

\begin{itemize}
\tightlist
\item
  Um oráculo de bit é \(U_f\ket{x,y}=\ket{x,y\oplus f(x)}\), unitário e auto-inverso.
\item
  Com o qubit extra em \(\ket{-}\), ele vira um oráculo de fase: \(\ket{x}\mapsto(-1)^{f(x)}\ket{x}\).
\item
  Ordem de bits: o mais significativo primeiro (big-endian); \(U_f\) tem tamanho exponencial se feito pela tabela e linear se feito por uma fórmula.
\item
  Complexidade de consultas conta chamadas ao oráculo; complexidade de circuito conta portas.
\end{itemize}

\begin{center}\rule{0.5\linewidth}{0.5pt}\end{center}

Fonte: texto original, com a estrutura do assunto seguindo B. Valiron, \emph{Introduction to Quantum Algorithms and Quantum Programming} (notas de curso, 2025), seções 5.1 a 5.4.

\subsection{Algoritmo de Deutsch-Jozsa}
Alguém promete que uma função \(f\) de \(n\) bits é de um de dois tipos e você precisa descobrir qual, perguntando o menos possível. Um computador clássico precisa de muitas perguntas no pior caso. Um quântico resolve com uma só.

\subsubsection{O problema}\label{o-problema}

Seja \(f:\{0,1\}^n\to\{0,1\}\) com a \textbf{promessa} de que \(f\) é:

\begin{itemize}
\tightlist
\item
  \textbf{constante}: \(f(x)\) é o mesmo para todo \(x\) (sempre \(0\) ou sempre \(1\)); ou
\item
  \textbf{balanceada}: \(f(x)=0\) para exatamente metade das \(2^n\) entradas e \(f(x)=1\) para a outra metade.
\end{itemize}

A tarefa é dizer qual dos dois casos vale. Com \(n=2\), a função de tabela \((f(00),f(01),f(10),f(11))=(0,1,1,0)\) é balanceada; \((1,1,1,1)\) é constante; \((0,0,0,1)\) viola a promessa e não nos interessa.

\textbf{Custo clássico.} Cada consulta dá um valor. Se as primeiras \(2^{n-1}\) respostas forem todas iguais, ainda pode ser constante ou balanceada (a metade restante poderia ser diferente), então só com \(2^{n-1}+1\) respostas iguais temos certeza de que é constante. No pior caso, são \(2^{n-1}+1\) consultas para uma resposta sempre correta.

Com \(n=3\), quantas consultas clássicas o pior caso exige?

\subsubsection{\texorpdfstring{Começando pequeno: \(n=1\) (Deutsch)}{Começando pequeno: n=1 (Deutsch)}}\label{comeuxe7ando-pequeno-n1-deutsch}

Aqui \(f\) tem só dois valores, \(f(0)\) e \(f(1)\). Constante significa \(f(0)=f(1)\); balanceada significa \(f(0)\neq f(1)\). Clássico: duas consultas (precisa ver os dois valores). Quântico, com uma:

\begin{enumerate}
\def\labelenumi{\arabic{enumi}.}
\tightlist
\item
  Comece com \(\ket{0}\ket{1}\) e aplique \(H\otimes H\): \(\ket{+}\ket{-}\).
\item
  Aplique \(U_f\). Pelo phase kickback (aula anterior), o primeiro qubit vira \(\tfrac{1}{\sqrt2}\big((-1)^{f(0)}\ket0+(-1)^{f(1)}\ket1\big)\) e o segundo continua em \(\ket{-}\).
\item
  Aplique \(H\) no primeiro qubit e meça.
\end{enumerate}

Veja o resultado em cada caso:

\begin{itemize}
\tightlist
\item
  \(f\) constante: o primeiro qubit é \(\pm\ket{+}\), e \(H\ket{+}=\ket0\). Mede \(0\).
\item
  \(f\) balanceada: é \(\pm\ket{-}\), e \(H\ket{-}=\ket1\). Mede \(1\).
\end{itemize}

A medida \(0\) diz ``constante'' e \(1\) diz ``balanceada'', com certeza, usando \textbf{uma} consulta.

Com \(f(0)=0\) e \(f(1)=1\) (função balanceada), o que o algoritmo conclui? Responda ``constante'' ou ``balanceada''.

\subsubsection{O caso geral}\label{o-caso-geral}

Antes da fórmula, veja o caso \(n=2\) com \(f=(0,1,1,0)\): depois do oráculo o registrador é \(\tfrac12(\ket{00}-\ket{01}-\ket{10}+\ket{11})=\ket-\ket-\), e \(H\otimes H\) leva isso a \(\ket{11}\), nunca a \(\ket{00}\). Agora o caso geral.

O circuito é o mesmo, com \(H\) no registrador inteiro de \(n\) qubits:

\begin{enumerate}
\def\labelenumi{\arabic{enumi}.}
\tightlist
\item
  Prepare \(\ket{0}^{\otimes n}\ket{1}\).
\item
  Aplique \(H^{\otimes n}\) nos \(n\) primeiros qubits e \(H\) no último: o registrador fica na superposição uniforme e o último em \(\ket-\).
\item
  Aplique \(U_f\) (que age como oráculo de fase no primeiro registrador).
\item
  Aplique \(H^{\otimes n}\) no primeiro registrador e \textbf{meça-o}.
\end{enumerate}

Para acompanhar as contas, precisamos do que \(H^{\otimes n}\) faz num estado da base:

\[
H^{\otimes n}\ket{x}=\frac{1}{\sqrt{2^n}}\sum_{z\in\{0,1\}^n}(-1)^{x\cdot z}\ket{z},\qquad x\cdot z=x_1z_1\oplus\cdots\oplus x_nz_n,
\]

onde \(x\cdot z\) é o \textbf{produto escalar módulo 2} (soma dos produtos bit a bit, mod 2). Isso vem de \(H\ket{b}=\tfrac1{\sqrt2}\sum_{c}(-1)^{bc}\ket c\), aplicado em cada qubit.

Depois dos passos 2 e 3, o primeiro registrador está em

\[
\frac{1}{\sqrt{2^n}}\sum_x(-1)^{f(x)}\ket{x}.
\]

Depois do passo 4, a amplitude do estado \(\ket{z}\) é

\[
\frac{1}{2^n}\sum_x(-1)^{f(x)+x\cdot z}.
\]

Só nos interessa \(z=0\cdots0\), onde \(x\cdot z=0\) para todo \(x\):

\[
\text{amplitude de }\ket{0\cdots0}=\frac{1}{2^n}\sum_x(-1)^{f(x)}.
\]

\begin{itemize}
\tightlist
\item
  \textbf{\(f\) constante}: todos os termos valem \(+1\) (ou todos \(-1\)), a soma é \(\pm2^n\) e a amplitude é \(\pm1\). Mede \(0\cdots0\) com probabilidade 1.
\item
  \textbf{\(f\) balanceada}: metade dos termos é \(+1\) e metade é \(-1\); a soma é \(0\). A probabilidade de medir \(0\cdots0\) é zero, então algum outro resultado aparece.
\end{itemize}

Portanto: se o resultado for \(0\cdots0\), \(f\) é constante; qualquer outro resultado, balanceada. Uma consulta, sempre correto.

\textbf{Exemplo com \(n=2\).} Seja \(f\) constante e igual a \(1\). A amplitude de \(\ket{00}\) é \(\tfrac14(-1-1-1-1)\). Quanto vale?

Agora a balanceada \(f(x_1x_2)=x_1\oplus x_2\), de tabela \((0,1,1,0)\). A soma \(\sum_x(-1)^{f(x)}=1-1-1+1=0\), então a probabilidade de medir \(00\) é:

\subsubsection{Quanto se ganha?}\label{quanto-se-ganha}

O algoritmo quântico usa \textbf{1} consulta; o clássico determinístico precisa de \(2^{n-1}+1\). Na prática, o clássico \textbf{probabilístico} com \(k\) consultas aleatórias erra com probabilidade de no máximo \(2^{1-k}\), sem depender de \(n\) (com \(k=20\), cerca de \(10^{-6}\)), e o erro é só de um lado. Então a vantagem real do DJ é sobre algoritmos \textbf{exatos}. Ainda assim, ele foi um dos primeiros exemplos de separação entre computação clássica e quântica em consultas.

Simule o algoritmo: dada a tabela de \(f\), calcule a probabilidade de medir \(0\cdots0\).

Nota: \texttt{f{[}x{]}} usa \(x\) como inteiro com \(x_1\) mais significativo, e \texttt{n} é \(\log_2\) de \texttt{len(f)}. O teste com \((0,0,0,1)\) não satisfaz a promessa; só calcule a fórmula.

Dica: monte \(H^{\otimes n}\) com \texttt{np.kron}, aplique o oráculo de fase \texttt{np.diag((-1.0)**np.array(f))} ao estado uniforme, aplique \(H^{\otimes n}\) de novo e pegue \(|\text{amplitude}[0]|^2\).

\subsubsection{Resumo}\label{resumo}

\begin{itemize}
\tightlist
\item
  Promessa: \(f\) constante ou balanceada. Clássico exato: \(2^{n-1}+1\) consultas.
\item
  Deutsch-Jozsa: \(H^{\otimes n}\), oráculo (com \(\ket-\) no qubit extra), \(H^{\otimes n}\), medir. Resultado \(0\cdots0\) significa constante.
\item
  A chave é que a amplitude de \(\ket{0\cdots0}\) é \(2^{-n}\sum_x(-1)^{f(x)}\).
\end{itemize}

\begin{center}\rule{0.5\linewidth}{0.5pt}\end{center}

Fonte: texto original; o assunto segue B. Valiron, \emph{Introduction to Quantum Algorithms and Quantum Programming} (notas de curso, 2025), seção 6.1.

\subsection{Algoritmo de Bernstein-Vazirani}
Agora a caixa-preta esconde um segredo: uma string de bits \(s\). Cada consulta devolve só \textbf{um bit} de informação sobre \(s\). Quantas consultas para descobrir o segredo inteiro?

\subsubsection{O problema}\label{o-problema}

Seja \(s\in\{0,1\}^n\) uma string secreta e

\[
f(x)=s\cdot x \bmod 2 = s_1x_1\oplus s_2x_2\oplus\cdots\oplus s_nx_n.
\]

Ou seja, \(f(x)\) é a \textbf{paridade} dos bits de \(x\) nas posições onde \(s\) vale \(1\). A tarefa é achar \(s\).

\textbf{Exemplo com \(n=3\) e \(s=101\).} Então \(f(x)=x_1\oplus x_3\). Para \(x=111\): \(1\oplus1=0\). Para \(x=100\): \(1\).

Qual é \(f(111)\) quando \(s=101\)?

\textbf{Custo clássico.} Consultar \(x=100\ldots0\) dá \(s_1\), consultar \(x=010\ldots0\) dá \(s_2\), e assim por diante: com \(n\) consultas lemos os \(n\) bits de \(s\). E não dá para fazer melhor, porque cada consulta traz só um bit e \(s\) tem \(n\) bits. Com \(n=8\), quantas consultas clássicas são necessárias?

\subsubsection{O algoritmo quântico: uma consulta}\label{o-algoritmo-quuxe2ntico-uma-consulta}

Surpresa: o circuito é \textbf{idêntico} ao de Deutsch-Jozsa.

\begin{enumerate}
\def\labelenumi{\arabic{enumi}.}
\tightlist
\item
  Prepare \(\ket{0}^{\otimes n}\ket{1}\) e aplique \(H\) em tudo.
\item
  Aplique \(U_f\) (com o qubit extra em \(\ket-\) ele age como o oráculo de fase \(\ket{x}\mapsto(-1)^{f(x)}\ket x\)).
\item
  Aplique \(H^{\otimes n}\) no primeiro registrador e meça.
\end{enumerate}

Vamos às contas. Como no DJ, antes do último \(H^{\otimes n}\) o registrador está em \(\tfrac{1}{\sqrt{2^n}}\sum_x(-1)^{s\cdot x}\ket x\). Depois do último, a amplitude de \(\ket z\) é

\[
\frac{1}{2^n}\sum_x(-1)^{s\cdot x+z\cdot x}=\frac{1}{2^n}\sum_x(-1)^{(s\oplus z)\cdot x}.
\]

Aqui usamos que \(s\cdot x\oplus z\cdot x=(s\oplus z)\cdot x\) (soma bit a bit). Agora, dois casos:

\begin{itemize}
\tightlist
\item
  Se \(z=s\), então \(s\oplus z=0\cdots0\), todos os termos valem \(+1\) e a soma é \(2^n\): amplitude \(1\).
\item
  Se \(z\ne s\), \(s\oplus z\) tem algum bit \(1\), e metade dos \(x\) dá \((s\oplus z)\cdot x=0\) e a outra metade dá \(1\). A soma é zero.
\end{itemize}

Logo o estado final é exatamente \(\ket{s}\) e a medida devolve \(s\) \textbf{com probabilidade 1}, com uma única consulta.

\textbf{Verificação com \(n=2\) e \(s=10\).} Então \(f(x_1x_2)=x_1\), de tabela \((0,0,1,1)\). O estado antes do último \(H^{\otimes2}\) é \(\tfrac12(\ket{00}+\ket{01}-\ket{10}-\ket{11})=\ket{-}\ket{+}\). Aplicar \(H\otimes H\) dá \(\ket{1}\ket{0}=\ket{10}=\ket s\). Funciona.

Para \(n=3\) e \(f\) de tabela \((0,1,1,0,1,0,0,1)\) (entradas \(x=000,001,\dots,111\)), qual string \(s\) o algoritmo devolve?

\subsubsection{Por que isso é interessante}\label{por-que-isso-uxe9-interessante}

DJ e BV usam o mesmo circuito, mas respondem perguntas diferentes: DJ distingue constante de balanceada (e o resultado \(0\cdots0\) ou não); BV lê \textbf{todas} as \(n\) componentes de \(s\) numa só consulta. A informação que cabe em \(n\) bits sai de uma chamada porque o oráculo age em todos os ramos da superposição e só a interferência do \(H^{\otimes n}\) final concentra a amplitude em \(\ket s\).

Note que isso é uma vantagem de \textbf{consultas} (\(1\) contra \(n\)), polinomial, não exponencial. Foi um passo importante para algoritmos posteriores, que usam a mesma estrutura de transformada de Hadamard para extrair uma propriedade escondida da função.

\subsubsection{Simulação}\label{simulauxe7uxe3o}

Implemente o algoritmo: dada a tabela de \(f(x)=s\cdot x \bmod 2\), devolva a string \(s\) encontrada, como texto de \(n\) bits.

Nota: \texttt{f{[}x{]}} usa \(x\) como inteiro com \(x_1\) mais significativo, e \texttt{n} é \(\log_2\) de \texttt{len(f)}.

Dica: reutilize sua simulação do Deutsch-Jozsa e, no fim, pegue o índice de maior amplitude com \texttt{np.argmax(np.abs(...))}, escrito em binário com \texttt{format(i,\ "0\%db"\ \%\ n)}.

\subsubsection{Resumo}\label{resumo}

\begin{itemize}
\tightlist
\item
  Problema: \(f(x)=s\cdot x\bmod2\), achar \(s\). Clássico: \(n\) consultas.
\item
  Quântico: o mesmo circuito do Deutsch-Jozsa; a amplitude de \(\ket z\) é \(2^{-n}\sum_x(-1)^{(s\oplus z)\cdot x}\), que vale \(1\) só para \(z=s\).
\item
  Uma consulta, sucesso com probabilidade 1.
\end{itemize}

\begin{center}\rule{0.5\linewidth}{0.5pt}\end{center}

Fonte: texto original; o assunto segue B. Valiron, \emph{Introduction to Quantum Algorithms and Quantum Programming} (notas de curso, 2025), seção 6.1.

\subsection{Busca de Grover e amplificação de amplitude}
Uma lista com \(N\) itens sem nenhuma ordem, e um único item é o que você procura. Sem estrutura, não há atalho clássico: no pior caso é preciso olhar quase todos. O algoritmo de Grover acha o item com cerca de \(\sqrt N\) consultas.

\subsubsection{O problema}\label{o-problema}

Seja \(f:\{0,1\}^n\to\{0,1\}\) com \(N=2^n\) entradas. Sabemos que existe um (ou poucos) \(w\) \textbf{marcado}, com \(f(w)=1\), e \(f(x)=0\) nos outros. A tarefa é encontrar \(w\).

Clássico: em média \(N/2\) consultas, e \(N-1\) no pior caso. Quântico (Grover): da ordem de \(\sqrt N\). Veja a escala: com \(N=10^6\), são cerca de \(10^6\) consultas contra cerca de \(785\).

\subsubsection{Ferramentas: duas reflexões}\label{ferramentas-duas-reflexuxf5es}

O algoritmo é feito de duas \textbf{reflexões}, e a ideia geométrica é: duas reflexões seguidas formam uma rotação.

\textbf{Estado de partida.} Aplicando \(H^{\otimes n}\) em \(\ket{0\cdots0}\) obtemos o estado uniforme

\[
\ket{\psi}=\frac{1}{\sqrt N}\sum_x\ket x.
\]

Medi-lo agora acharia \(w\) com probabilidade \(1/N\), o mesmo que chutar.

\textbf{Reflexão 1: o oráculo de fase.} Com o oráculo de fase da primeira aula, \(O_f\ket x=(-1)^{f(x)}\ket x\), só \(\ket w\) troca de sinal:

\[
O_f=I-2\ket w\bra w.
\]

Isso é uma reflexão: mantém o que é ortogonal a \(\ket w\) e inverte a componente em \(\ket w\).

\textbf{Reflexão 2: a difusão.} Defina

\[
D=2\ket\psi\bra\psi-I,
\]

uma reflexão em torno de \(\ket\psi\) (mantém o que é paralelo a \(\ket\psi\), inverte o resto). Na prática, \(D=H^{\otimes n}(2\ket{0}\bra{0}-I)H^{\otimes n}\), e \(2\ket0\bra0-I\) é ``trocar o sinal de todos os estados exceto \(\ket{0\cdots0}\)'', o que é equivalente, a menos de fase global, a uma porta \(Z\) multicontrolada cercada de portas \(X\) em todos os qubits. Essa fase global \(-1\) passa a importar se \(G\) for usada de forma controlada (por exemplo, em contagem de amplitude).

O efeito de \(D\) em um estado \(\sum a_x\ket x\) com média \(\bar a=\tfrac1N\sum a_x\) é \(a_x\mapsto 2\bar a-a_x\): cada amplitude é \textbf{refletida em torno da média}.

\textbf{Um passo de Grover} é \(G=D\,O_f\).

\subsubsection{\texorpdfstring{Exemplo com \(N=4\) (\(n=2\)), item marcado \(w=11\)}{Exemplo com N=4 (n=2), item marcado w=11}}\label{exemplo-com-n4-n2-item-marcado-w11}

Começamos com \(\ket\psi=\tfrac12(1,1,1,1)\) (amplitudes de \(\ket{00},\ket{01},\ket{10},\ket{11}\)). Pense nas amplitudes como alturas de barras:

\begin{itemize}
\tightlist
\item
  Inicial: quatro barras de altura \(\tfrac12\).
\item
  Depois do oráculo: troca o sinal de \(\ket{11}\), e a barra dele vira \(-\tfrac12\) (fica abaixo do zero). Barras: \(\tfrac12,\tfrac12,\tfrac12,-\tfrac12\).
\item
  A \textbf{média} das amplitudes caiu para \(\bar a=\tfrac14\).
\item
  Difusão: cada barra é refletida em torno da média, \(a\mapsto2\bar a-a\). As três primeiras ficam em \(2\cdot\tfrac14-\tfrac12=0\); a última em \(2\cdot\tfrac14+\tfrac12=1\).
\end{itemize}

Resultado: \((0,0,0,1)=\ket{11}\). Com \textbf{uma} iteração, achamos \(w\) com probabilidade 1; um clássico precisaria de até 3 consultas.

O efeito geral é o mesmo: a difusão \(D\) reflete cada amplitude em torno da média. Como o oráculo puxa a média para baixo, a barra marcada, que ficou negativa, vira a mais alta depois da reflexão.

\subsubsection{Visão geométrica}\label{visuxe3o-geomuxe9trica}

Todo o jogo acontece num \textbf{plano} de dimensão 2. Defina

\[
\ket{w^\perp}=\frac{1}{\sqrt{N-1}}\sum_{x\ne w}\ket x,
\]

o estado uniforme sobre os itens não marcados. Então \(\ket w\) e \(\ket{w^\perp}\) são ortonormais e

\[
\ket\psi=\sin\theta\,\ket w+\cos\theta\,\ket{w^\perp},\qquad \sin\theta=\frac{1}{\sqrt N}.
\]

Isto é, \(\ket\psi\) forma o ângulo \(\theta\) com \(\ket{w^\perp}\) (bem perto, quando \(N\) é grande: \(\theta\approx1/\sqrt N\)).

\begin{itemize}
\tightlist
\item
  \(O_f\) reflete em torno do eixo \(\ket{w^\perp}\): o ângulo medido a partir de \(\ket{w^\perp}\) troca de sinal, de \(\theta\) para \(-\theta\).
\item
  \(D\) reflete em torno de \(\ket\psi\) (ângulo \(\theta\)): o ângulo \(-\theta\) vira \(2\theta-(-\theta)=3\theta\).
\end{itemize}

Cada iteração gira o estado de \(2\theta\) em direção a \(\ket w\). Depois de \(k\) iterações:

\[
G^k\ket\psi=\sin\big((2k+1)\theta\big)\ket w+\cos\big((2k+1)\theta\big)\ket{w^\perp},
\]

e a probabilidade de medir \(w\) é \(\sin^2\!\big((2k+1)\theta\big)\).

Confira com \(N=4\): \(\sin\theta=\tfrac12\), \(\theta=30^\circ\), e para \(k=1\), \((2k+1)\theta=90^\circ\), \(\sin^2=1\). Exatamente o que vimos.

Para \(N=8\) e \(k=1\) a amplitude do item marcado fica em \(\tfrac{5}{4\sqrt2}\approx0{,}884\): não chega a 1, mas é bem melhor que \(1/\sqrt8\). Pela fórmula, \(\sin\theta=1/\sqrt8\), \(\theta\approx20{,}7^\circ\), \(3\theta\approx62{,}1^\circ\). Qual é a probabilidade de sucesso, \(\sin^2 3\theta\)? (Dica: \(\sin3\theta=3\sin\theta-4\sin^3\theta\).)

Agora à mão: reflita em torno da média as amplitudes \((\tfrac12,\tfrac12,\tfrac12,-\tfrac12)\) (média \(\tfrac14\)). Qual é a nova amplitude do item que estava em \(-\tfrac12\)?

Com \(N=4\), uma iteração dá probabilidade \(1\). E se aplicarmos uma segunda? Calcule \(\sin^2(5\theta)\) com \(\theta=30^\circ\).

\subsubsection{Quantas iterações?}\label{quantas-iterauxe7uxf5es}

Queremos \((2k+1)\theta\approx\pi/2\), isto é,

\[
k\approx\frac{\pi}{4\theta}-\frac12,\qquad \sin\theta=\sqrt{M/N},
\]

arredondado para o inteiro mais próximo, onde \(M\) é o número de itens marcados (\(M=1\) no caso básico). Para \(N\) grande, isso é cerca de \(\tfrac\pi4\sqrt{N/M}\): da ordem de \(\sqrt N\). Mais itens marcados, menos iterações.

Com \(N=1024\) e \(M=1\) (\(\theta=\arcsin(1/32)\approx0{,}03126\)), quantas iterações?

\textbf{Atenção: não vale ``mais é melhor''.} A rotação não para sozinha. Se passar de \(\pi/2\), o estado se afasta de \(\ket w\) e a probabilidade cai, oscilando. Por isso é preciso conhecer \(M\); se não o conhecemos, usa-se um número aleatório de iterações (algoritmo BBHT) ou contagem quântica.

Com \(N=64\) e \(M=4\) itens marcados (\(\sin\theta=1/4\), \(\theta\approx0{,}2527\) rad), quantas iterações?

\subsubsection{Simulação}\label{simulauxe7uxe3o}

Simule Grover com NumPy: monte o estado uniforme, aplique \(k\) vezes o oráculo (\(-1\) em \(w\)) e a difusão (refletir em torno da média), e devolva a probabilidade de medir \(w\).

Dica: com \texttt{psi\ =\ np.ones(N)\ /\ np.sqrt(N)}, o oráculo é \texttt{psi{[}marcado{]}\ *=\ -1} e a difusão é \texttt{psi\ =\ 2\ *\ psi.mean()\ -\ psi}.

\subsubsection{Para ir além: amplitude amplification (opcional)}\label{para-ir-aluxe9m-amplitude-amplification-opcional}

Esta seção generaliza o Grover e pode ser pulada numa primeira leitura.

O método não precisa começar no estado uniforme. Suponha um algoritmo quântico \(\mathcal A\) (um circuito unitário) que prepara

\[
\mathcal A\ket{0\cdots0}=\sin\theta\,\ket{\text{bom}}+\cos\theta\,\ket{\text{ruim}},
\]

onde ``bom'' e ``ruim'' são estados normalizados e ortogonais, e \(p=\sin^2\theta\) é a probabilidade de medir um resultado bom. Uma tentativa direta acerta com probabilidade \(p\); repetir dá em média \(1/p\) tentativas.

O mesmo truque das duas reflexões funciona:

\begin{itemize}
\tightlist
\item
  reflexão 1: troca o sinal da parte ``boa'' (o oráculo, que reconhece resultados bons);
\item
  reflexão 2: reflexão em torno do estado inicial \(\mathcal A\ket{0\cdots0}\), feita por \(\mathcal A\,(2\ket0\bra0-I)\,\mathcal A^{-1}\).
\end{itemize}

Cada rodada gira de \(2\theta\) e, depois de \(k\) rodadas, a probabilidade de sucesso é \(\sin^2((2k+1)\theta)\). Como \(\theta\approx\sqrt p\), bastam cerca de \(\tfrac\pi4\sqrt{1/p}\) rodadas em vez de \(1/p\): um ganho \textbf{quadrático}. O Grover é o caso em que \(\mathcal A=H^{\otimes n}\) e \(p=1/N\).

Se um algoritmo acerta com probabilidade \(p=1/100\), quantas rodadas de amplificação (mesma fórmula, com \(\sin\theta=\sqrt p\))?

Esse resultado é \textbf{ótimo}: nenhum algoritmo quântico acha um item marcado numa lista não estruturada com bem menos que \(\sqrt N\) consultas (Bennett, Bernstein, Brassard e Vazirani, 1997; Zalka, 1999, mostra que a constante \(\pi/4\) também é a melhor possível). Por isso Grover dá um ganho quadrático, não exponencial.

\subsubsection{Resumo}\label{resumo}

\begin{itemize}
\tightlist
\item
  Problema: achar \(w\) com \(f(w)=1\) entre \(N=2^n\) entradas, sem estrutura.
\item
  Passo de Grover: oráculo de fase (reflexão que inverte \(\ket w\)) e difusão (reflexão em torno de \(\ket\psi\)); juntas, uma rotação de \(2\theta\), com \(\sin\theta=1/\sqrt N\).
\item
  Depois de \(k\) passos, \(P(w)=\sin^2((2k+1)\theta)\); o melhor \(k\) é o inteiro mais próximo de \(\tfrac{\pi}{4\theta}-\tfrac12\) (cerca de \(\tfrac\pi4\sqrt N\)).
\item
  Opcional: a amplitude amplification generaliza, e qualquer algoritmo com sucesso \(p\) vira \(\sim\sqrt{1/p}\) rodadas.
\end{itemize}

\begin{center}\rule{0.5\linewidth}{0.5pt}\end{center}

Fonte: texto original; o assunto segue B. Valiron, \emph{Introduction to Quantum Algorithms and Quantum Programming} (notas de curso, 2025), seção 5.4.

\subsection{Transformada de Fourier quântica}
Quase todo algoritmo quântico que dá ganho exponencial esconde, no fim, a mesma pergunta: ``que \textbf{frequência} tem esse padrão?''. A transformada de Fourier responde isso, e a versão quântica faz a conta com muito poucas portas. Aqui você vai construí-la primeiro com números pequenos, depois em geral.

\subsubsection{A transformada de Fourier discreta (DFT)}\label{a-transformada-de-fourier-discreta-dft}

Um sinal com \(N\) amostras é um vetor \(x = (x_0,\dots,x_{N-1})\). A DFT troca ``valor em cada instante'' por ``quanto de cada frequência''. Ela usa as \textbf{raízes da unidade}

\[
\omega = e^{2\pi i/N},
\]

que são \(N\) pontos igualmente espaçados no círculo unitário (\(\omega^N = 1\)). A saída é

\[
y_k = \frac{1}{\sqrt N}\sum_{j=0}^{N-1} \omega^{jk}\, x_j .
\]

Exemplo com \(N=4\): aqui \(\omega = e^{2\pi i/4} = i\), e as potências são \(\omega^0=1,\ \omega^1=i,\ \omega^2=-1,\ \omega^3=-i,\ \omega^4=1,\dots\) (o expoente só importa módulo 4). A matriz \(F_{kj}=\omega^{jk}/\sqrt N\) fica

\[
F_4 = \frac12\begin{pmatrix} 1&1&1&1\\ 1&i&-1&-i\\ 1&-1&1&-1\\ 1&-i&-1&i \end{pmatrix}.
\]

Para \(N=2\), \(\omega=-1\) e \(F_2\) é a matriz de Hadamard. Calculada direto, a DFT custa cerca de \(N^2\) multiplicações. A \textbf{FFT} (transformada rápida) baixa isso para cerca de \(N\log N\) separando as posições pares e ímpares: uma DFT de tamanho \(N\) vira duas de tamanho \(N/2\), e repete-se. Mesmo assim, cresce com \(N\).

Qual é o ângulo (em graus, entre 0 e 360) de \(\omega^{jk}\) para \(N=8\), \(j=3\) e \(k=5\)? (Dica: \(jk=15\) e só vale \(jk\) módulo 8.)

\subsubsection{A QFT: a mesma conta nas amplitudes}\label{a-qft-a-mesma-conta-nas-amplitudes}

Com \(n\) qubits temos \(N=2^n\) estados da base. A \textbf{transformada de Fourier quântica} (QFT) é a matriz unitária \(F\) acima, aplicada ao vetor de amplitudes:

\[
\mathrm{QFT}\ket{j} = \frac{1}{\sqrt N}\sum_{k=0}^{N-1} \omega^{jk}\ket{k}.
\]

Ela é unitária (as colunas são ortonormais), então é uma porta legítima. Para \(n=2\), a primeira coluna de \(F_4\) diz \(\mathrm{QFT}\ket{00} = \tfrac12(\ket{00}+\ket{01}+\ket{10}+\ket{11})\), a superposição uniforme, igual a \(H\otimes H\ket{00}\). Qual é a amplitude de cada estado nesse caso?

Já \(\mathrm{QFT}\ket{01}\) é a segunda coluna: \(\tfrac12(\ket{00}+i\ket{01}-\ket{10}-i\ket{11})\). Todas as amplitudes têm o mesmo módulo, mas as \textbf{fases giram} com uma ``frequência'' que depende de \(j\). Essa é a mensagem central: a QFT leva um número \(j\) guardado na base para uma onda de frequência \(j\) guardada nas fases.

Atenção: a QFT \textbf{não} faz a DFT ``de graça'' sobre dados clássicos. Ela transforma amplitudes, que não podemos ler todas. Seu valor está em usar o resultado dentro de outro algoritmo.

\subsubsection{O circuito para n = 2, passo a passo}\label{o-circuito-para-n-2-passo-a-passo}

Para \(n=2\) escrevemos \(j=2j_1+j_2\) (então \(j_1\) é o bit mais significativo). Precisamos de uma porta de fase \(R_m=\begin{pmatrix}1&0\\0&e^{2\pi i/2^m}\end{pmatrix}\); aqui basta \(R_2=S\) (fase \(e^{2\pi i/4}=i\)). O circuito:

\begin{enumerate}
\def\labelenumi{\arabic{enumi}.}
\tightlist
\item
  \(H\) no qubit 1.
\item
  \(S\) no qubit 1, controlada pelo qubit 2.
\item
  \(H\) no qubit 2.
\item
  SWAP entre os dois qubits.
\end{enumerate}

Veja o que acontece com a entrada \(j=1\) (isto é, \(j_1=0\), \(j_2=1\), estado \(\ket{01}\)):

\begin{itemize}
\tightlist
\item
  \(H\) no qubit 1: \(\tfrac1{\sqrt2}(\ket0+\ket1)\), pois \(j_1=0\).
\item
  \(S\) controlada: o qubit 2 vale 1, então o \(\ket1\) do qubit 1 ganha a fase \(i\): \(\tfrac1{\sqrt2}(\ket0+i\ket1)\).
\item
  \(H\) no qubit 2 (que está em \(\ket1\)): \(\tfrac1{\sqrt2}(\ket0-\ket1)\).
\item
  Antes do SWAP, o estado é \(\tfrac1{\sqrt2}(\ket0+i\ket1)\otimes\tfrac1{\sqrt2}(\ket0-\ket1)\). Depois do SWAP: \(\tfrac12(\ket0-\ket1)(\ket0+i\ket1)=\tfrac12(\ket{00}+i\ket{01}-\ket{10}-i\ket{11})\), a segunda coluna de \(F_4\).
\end{itemize}

A saída é um \textbf{produto} de dois qubits (sem emaranhamento), e cada qubit recebeu uma fase que depende de alguns bits de \(j\). Qual é o ângulo (em graus, entre 0 e 360) da amplitude de \(\ket{01}\) no estado \(\mathrm{QFT}\ket{10}\) (que é \(j=2\), a terceira coluna de \(F_4\))?

\subsubsection{Fração binária e o caso geral}\label{frauxe7uxe3o-binuxe1ria-e-o-caso-geral}

Para generalizar usamos a \textbf{fração binária}

\[
0.j_1j_2\cdots j_m = \frac{j_1}{2}+\frac{j_2}{4}+\cdots+\frac{j_m}{2^m}.
\]

Exemplo: \(0.101 = \tfrac12+\tfrac18 = 0{,}625\). Para \(n=3\) e \(j=j_1j_2j_3\), o resultado da QFT é, por extenso,

\[
\mathrm{QFT}\ket{j_1j_2j_3}=\tfrac1{\sqrt8}\big(\ket0+e^{2\pi i\,0.j_3}\ket1\big)\big(\ket0+e^{2\pi i\,0.j_2j_3}\ket1\big)\big(\ket0+e^{2\pi i\,0.j_1j_2j_3}\ket1\big).
\]

Para \(n=2\) são os dois últimos fatores, que você viu acima (com \(j=1\): \(0.j_2=\tfrac12\) e \(0.j_1j_2=\tfrac14\)). Em geral, o qubit \(l\) da saída leva a fase \(e^{2\pi i\,0.j_{n-l+1}\cdots j_n}\).

Receita do circuito, qubit a qubit (com \(R_m\) de fase \(e^{2\pi i/2^m}\); casos conhecidos: \(R_1=Z\), \(R_2=S\) e \(R_3=T\), com fase de 45°):

\begin{enumerate}
\def\labelenumi{\arabic{enumi}.}
\tightlist
\item
  No qubit 1: \(H\), depois \(R_2\) controlada pelo qubit 2, \(R_3\) controlada pelo qubit 3, \ldots, \(R_n\) controlada pelo qubit \(n\).
\item
  No qubit 2: \(H\), depois \(R_2\) controlada pelo qubit 3, \ldots, \(R_{n-1}\) controlada pelo qubit \(n\).
\item
  Continue até o qubit \(n\), que só recebe \(H\).
\item
  No fim, \textbf{inverta a ordem} dos qubits com SWAPs: sem isso, o qubit 1 teria a fase \(0.j_1\cdots j_n\), mas a fórmula manda a fase \(0.j_n\) para ele.
\end{enumerate}

Quantas portas controladas \(R_m\) o circuito tem para \(n\) qubits? O qubit 1 usa \(n-1\), o 2 usa \(n-2\), e assim por diante: \((n-1)+(n-2)+\cdots+1 = n(n-1)/2\). Quantas para \(n=5\)?

No total são \(n\) portas \(H\), \(n(n-1)/2\) fases controladas e cerca de \(n/2\) SWAPs: da ordem de \(n^2\) portas, contra os \(N\log N = 2^n\,n\) passos da FFT. Esse é o ganho exponencial, com a ressalva de que a saída é um estado quântico.

Se aplicarmos a QFT e depois a QFT inversa a \(\ket j\), que estado obtemos?

\subsubsection{QFT inversa}\label{qft-inversa}

Como \(F\) é unitária, a inversa é a adjunta (conjugada transposta). Como \(F\) é simétrica, basta conjugar: \(F^{-1}_{kj}=\omega^{-jk}/\sqrt N\). No circuito: percorra as portas \textbf{de trás para a frente} e troque cada \(R_m\) por \(R_m^\dagger\) (fase \(e^{-2\pi i/2^m}\)). A QFT inversa transforma ``onda de frequência \(j\)'' de volta em \(\ket{j}\), e é por isso que ela aparece no fim da estimação de fase.

\subsection{Estimação de fase}
Uma porta quântica \(U\) tem autovetores, e cada autovetor só ganha uma fase quando passa por ela. Descobrir essa fase é como ler a ``velocidade de rotação'' escondida na porta. A \textbf{estimação de fase} faz isso, e é o motor por trás do algoritmo de Shor, do HHL e de simulações químicas.

\subsubsection{O problema}\label{o-problema}

Seja \(U\) unitária e \(\ket{u}\) um autovetor dela:

\[
U\ket{u} = e^{2\pi i\varphi}\ket{u}, \qquad 0\le\varphi<1.
\]

Como \(U\) é unitária, o autovalor tem módulo 1 e sempre dá para escrevê-lo assim; \(\varphi\) é a \textbf{fase} (em voltas completas: \(\varphi=\tfrac14\) é um quarto de volta, 90°). Temos acesso ao circuito de \(U\) (e a versões controladas dele) e a uma cópia de \(\ket{u}\). Queremos \(\varphi\) com \(t\) bits de precisão.

Exemplo: \(U=S\) e \(\ket{u}=\ket1\). Como \(S\ket1=i\ket1 = e^{2\pi i\cdot\frac14}\ket1\), a resposta é \(\varphi=\tfrac14=0.01_2\). Qual é \(\varphi\) para \(U=T\) e \(\ket{u}=\ket1\) (use \(T\ket1=e^{i\pi/4}\ket1\))?

\subsubsection{Exemplo com 2 qubits de registro, antes da regra geral}\label{exemplo-com-2-qubits-de-registro-antes-da-regra-geral}

Use \(U=S\), \(\ket{u}=\ket1\) e um registro de \(t=2\) qubits, mais o qubit-alvo que guarda \(\ket u\). O \textbf{registro} é o grupo de \(t\) qubits onde a fase vai aparecer; o \textbf{alvo} é o qubit (ou grupo) que guarda \(\ket u\).

\begin{enumerate}
\def\labelenumi{\arabic{enumi}.}
\tightlist
\item
  Ponha o registro em \(H\otimes H\ket{00}=\tfrac12(\ket{00}+\ket{01}+\ket{10}+\ket{11})\).
\item
  O qubit menos significativo do registro controla \(U^1=S\); o mais significativo controla \(U^2=S^2=Z\). Cada vez que o controle vale 1, o alvo \(\ket1\) devolve a fase \(i\) (para \(S\)) ou \(-1\) (para \(S^2\)). A fase \textbf{fica no registro} (esse é o efeito de ``kickback'' dos oráculos): o termo \(\ket{x}\) ganha \(i^{x}=e^{2\pi i\,x/4}\), com \(x\) o número escrito no registro.
\item
  Resultado: \(\tfrac12(\ket{00}+i\ket{01}-\ket{10}-i\ket{11})\), que é exatamente a segunda coluna de \(F_4\), ou seja, \(\mathrm{QFT}\ket{01}\).
\item
  Aplicando a QFT inversa, voltamos a \(\ket{01}\). Medindo, lemos \(y=1\), e \(\varphi = y/2^t = 1/4\). Acertamos.
\end{enumerate}

\subsubsection{A regra geral}\label{a-regra-geral}

Com \(t\) qubits no registro:

\begin{enumerate}
\def\labelenumi{\arabic{enumi}.}
\tightlist
\item
  Prepare o registro em \(\ket{0}^{\otimes t}\) e o alvo em \(\ket{u}\). Aplique \(H\) em todo o registro: \(\tfrac1{\sqrt{2^t}}\sum_{x}\ket{x}\ket{u}\).
\item
  Para cada \(k=0,\dots,t-1\), aplique \(U^{2^k}\) ao alvo, controlada pelo qubit do registro de peso \(2^k\). Como \(U^{2^k}\ket u=e^{2\pi i\varphi 2^k}\ket u\), o termo \(\ket x\), com \(x=\sum_k x_k2^k\), recebe a fase \(e^{2\pi i\varphi x}\):
  \[
  \frac{1}{\sqrt{2^t}}\sum_{x=0}^{2^t-1} e^{2\pi i\,\varphi\, x}\ket{x}\ \ket{u}.
  \]
\item
  Aplique a QFT inversa no registro e meça.
\end{enumerate}

Por que dá certo: compare com \(\mathrm{QFT}\ket y=\tfrac1{\sqrt{2^t}}\sum_x e^{2\pi i\,xy/2^t}\ket x\). Se \(\varphi=y/2^t\) exatamente (um número com \(t\) bits binários), a soma acima é \(\mathrm{QFT}\ket y\), e a QFT inversa entrega \(\ket y\) com certeza.

Repare no custo: se só soubermos repetir \(U\), o qubit de peso \(2^k\) pede \(2^k\) aplicações, e o total é \(1+2+\cdots+2^{t-1}=2^t-1\). Cada qubit extra no registro \textbf{dobra} o custo. (No Shor, calcula-se \(U^{2^k}\) direto, sem repetir.) Se \(t\) sobe de 5 para 6, por qual fator, aproximadamente, o número total de aplicações de \(U\) é multiplicado?

Um segundo exemplo, agora com \(t=4\): se \(\varphi=0.1011_2=\tfrac12+\tfrac18+\tfrac1{16}=\tfrac{11}{16}\), qual inteiro \(y\) o registro mostra?

E se \(U\ket u=-\ket u\) (autovalor \(-1=e^{2\pi i\cdot\frac12}\)), qual é \(\varphi\)?

\subsubsection{\texorpdfstring{E se \(\varphi\) não cabe em \(t\) bits?}{E se \textbackslash varphi não cabe em t bits?}}\label{e-se-varphi-nuxe3o-cabe-em-t-bits}

Exemplo: \(U=T\) e \(\ket u=\ket1\) dão \(\varphi=\tfrac18=0.001_2\), que precisa de 3 bits. Com apenas \(t=2\) qubits, o melhor que o registro pode escrever é \(y/4\), com \(y=0,1,2,3\) (ou seja, \(0,\ 0{,}25,\ 0{,}5,\ 0{,}75\)). Fazendo a conta (a fórmula está logo abaixo), as probabilidades de medir \(y=0,1,2,3\) são

\[
0{,}427,\quad 0{,}427,\quad 0{,}073,\quad 0{,}073 .
\]

Os dois valores mais próximos de \(\tfrac18\) (\(0\) e \(\tfrac14\)) são os mais prováveis, mas o resultado é \textbf{aleatório}; com \(t=3\) teríamos certeza. Em geral, a probabilidade de medir \(y\) é

\[
P(y)=\left|\frac{1}{2^t}\sum_{x=0}^{2^t-1}e^{2\pi i\,x\,(\varphi-y/2^t)}\right|^2 ,
\]

que vale 1 quando \(\varphi=y/2^t\) exatamente e fica pequena quando \(\varphi\) está longe de \(y/2^t\). Em resumo: mede-se, com boa probabilidade, uma fração de \(t\) bits \textbf{perto} de \(\varphi\), e cada qubit extra dobra a resolução. (Para quem quiser os números exatos: a fração de \(t\) bits mais próxima de \(\varphi\) sai com probabilidade de pelo menos \(4/\pi^2\approx0{,}405\), e para obter \(\varphi\) com erro menor que \(2^{-n}\) com probabilidade pelo menos \(1-\epsilon\) usa-se \(t=n+\lceil\log_2(2+\tfrac1{2\epsilon})\rceil\) qubits; veja a fonte.)

Confira no computador:

Quantos qubits de registro são necessários para representar \(\varphi=\tfrac18\) (da porta \(T\)) sem erro?

\subsubsection{Vários autoestados}\label{vuxe1rios-autoestados}

E se não tivermos um autovetor, mas um estado qualquer? Como os autovetores de \(U\) formam uma base, \(\ket\psi=\sum_j c_j\ket{u_j}\) com \(U\ket{u_j}=e^{2\pi i\varphi_j}\ket{u_j}\). Por linearidade, o circuito leva

\[
\ket\psi\ket{0}\ \longmapsto\ \sum_j c_j\,\ket{\tilde\varphi_j}\ket{u_j},
\]

onde \(\ket{\tilde\varphi_j}\) é o registro com a estimativa de \(\varphi_j\). Medindo o registro, obtemos \(\varphi_j\) com probabilidade \(|c_j|^2\), e o alvo \textbf{colapsa} em \(\ket{u_j}\). Isto é útil: no Shor, não sabemos construir os autoestados, mas sabemos construir uma soma deles.

Exemplo: \(U=T\) e \(\ket\psi=\tfrac{\sqrt3}{2}\ket0+\tfrac12\ket1\). \(T\ket0=\ket0\) (\(\varphi=0\)) e \(T\ket1\) tem \(\varphi=\tfrac18\). Qual é a probabilidade de medir \(\varphi=\tfrac18\) (com \(t\ge3\))?

Qual bloco do circuito, aplicado ao registro antes da medida, converte a fase codificada em um número binário legível?

\begin{quote}
Para se aprofundar: B. Valiron, \emph{Introduction to Quantum Algorithms and Quantum Programming} (notas de curso, 2025), seção 5.6. As probabilidades de acerto citadas são o resultado padrão de M. Nielsen e I. Chuang, \emph{Quantum Computation and Quantum Information}, cap. 5.
\end{quote}

\subsection{Algoritmo de Shor}
Dado um número composto \(N\), como achar um divisor? Não se conhece método clássico rápido (os melhores levam tempo sub-exponencial no número de dígitos, isto é, bem pior que polinomial), e o \textbf{RSA}, sistema de criptografia usado na internet, depende dessa dificuldade. O algoritmo de Shor fatora em tempo polinomial num computador quântico. A ideia central é que \textbf{fatorar se reduz a achar um período}, e achar períodos é exatamente o que a estimação de fase faz bem.

\subsubsection{Aritmética modular em um minuto}\label{aritmuxe9tica-modular-em-um-minuto}

Escrevemos \(a \bmod N\) para o resto da divisão de \(a\) por \(N\). Exemplo: \(49 \bmod 15 = 4\), pois \(49=3\cdot15+4\). Dizemos que \(a\equiv b \pmod N\) se têm o mesmo resto. O \textbf{mdc} (máximo divisor comum) de dois números é o maior inteiro que divide ambos: \(\gcd(12,18)=6\). O algoritmo de Euclides o calcula rápido: \(\gcd(a,b)=\gcd(b,\ a \bmod b)\), parando quando o resto é 0.

\subsubsection{\texorpdfstring{Do período aos fatores (parte clássica), com \(N=15\)}{Do período aos fatores (parte clássica), com N=15}}\label{do-peruxedodo-aos-fatores-parte-cluxe1ssica-com-n15}

Escolha um número \(a\) sem fator comum com \(N\), digamos \(a=7\). Calcule as potências de \(7\) módulo 15:

\begin{center}\begin{tabular}{llllllllll}
\hline
\(x\) & 0 & 1 & 2 & 3 & 4 & 5 & 6 & 7 & 8 \\
\hline

\hline

\(7^x \bmod 15\) & 1 & 7 & 4 & 13 & 1 & 7 & 4 & 13 & 1 \\
\end{tabular}\end{center}

A sequência é \textbf{periódica}. O menor \(r>0\) com \(a^r\equiv1 \pmod N\) é a \textbf{ordem} (ou período) de \(a\) módulo \(N\). Aqui \(r=4\).

Qual é o período de \(2\) módulo \(15\)? (Calcule \(2,4,8,\dots\) até voltar a 1.)

Como \(r=4\) é par, podemos olhar o ``meio do ciclo'': \(7^{r/2}=7^2=49\equiv4 \pmod{15}\). Então \(7^4-1 = (7^2-1)(7^2+1)=48\cdot50\), e \(48\cdot 50 = 2400 = 160\cdot15\) é múltiplo de 15. Os fatores \(48\) e \(50\) não são múltiplos de 15 (cada um carrega só uma parte: \(48\) é múltiplo de \(3\) e \(50\) é múltiplo de \(5\)). Os mdc's separam as partes:

\[
\gcd(7^2-1,\,15)=\gcd(48,15)=3,\qquad \gcd(7^2+1,\,15)=\gcd(50,15)=5,
\]

e \(3\cdot5=15\). O mesmo vale em geral. Se \(r\) é par, então \(a^r-1=(a^{r/2}-1)(a^{r/2}+1)\) é múltiplo de \(N\). Se além disso \(a^{r/2}\not\equiv-1\pmod N\) (o caso \(a^{r/2}\equiv1\) é impossível, pois \(r\) é o menor expoente), nenhum fator sozinho é múltiplo de \(N\) e os dois mdc's dão divisores de verdade:

\[
p=\gcd\!\big(a^{r/2}-1,\ N\big),\qquad q=\gcd\!\big(a^{r/2}+1,\ N\big).
\]

Quanto vale \(\gcd(a^{r/2}+1,15)\) com \(a=2\) e \(r=4\) (então \(a^{r/2}=4\))?

Quando falha? Se \(r\) for ímpar, ou se \(a^{r/2}\equiv-1 \pmod N\) (por exemplo \(a=14\), \(N=15\): \(r=2\) e \(14\equiv-1\)). Basta sortear outro \(a\); para \(N\) ímpar e não potência de primo, cada sorteio funciona com probabilidade de pelo menos \(1/2\). Se o sorteio já der \(\gcd(a,N)>1\), achamos um divisor sem fazer nada.

Em resumo, a parte clássica é rápida. O difícil é achar \(r\): para \(N\) de \(n\) bits, o \textbf{valor} de \(r\) pode chegar perto de \(N\) (exponencial em \(n\)), e calcular a tabela acima seria impraticável.

Implemente a parte clássica:

\subsubsection{\texorpdfstring{A parte quântica, 1: a porta \(U_a\)}{A parte quântica, 1: a porta U\_a}}\label{a-parte-quuxe2ntica-1-a-porta-u_a}

Seja \(n\) o número de bits de \(N\) (\(n=4\) para \(N=15\)). Defina a unitária \(U_a\) nos estados \(\ket y\) com \(y<N\):
\[
U_a\ket{y} = \ket{a\,y \bmod N}.
\]
Para \(N=15\) e \(a=7\): \(U_7\ket1=\ket7\), \(U_7\ket7=\ket{49\bmod15}=\ket4\), \(U_7\ket4=\ket{28\bmod15}=\ket{13}\) e \(U_7\ket{13}=\ket{91\bmod15}=\ket1\). A porta apenas \textbf{roda o ciclo} \(1\to7\to4\to13\to1\). É uma permutação (porque \(\gcd(a,N)=1\)), logo unitária. Ela faz a exponenciação modular: \(U_a^{x}\ket1=\ket{a^x\bmod N}\). Em circuito, calcula-se \(U_a^{2^k}\) diretamente como multiplicação por \(a^{2^k}\bmod N\) (obtido por quadrados sucessivos, fácil classicamente), então cada potência custa só um circuito de multiplicação modular de tamanho polinomial em \(n\).

\subsubsection{A parte quântica, 2: os autoestados}\label{a-parte-quuxe2ntica-2-os-autoestados}

Um ciclo tem ``ondas'' como autovetores. Para \(N=15\), \(a=7\), os quatro estados (escreva os coeficientes na ordem \(\ket1,\ket7,\ket4,\ket{13}\)):

\[
\ket{u_0}=\tfrac12(1,\,1,\,1,\,1),\quad
\ket{u_1}=\tfrac12(1,\,-i,\,-1,\,i),\quad
\ket{u_2}=\tfrac12(1,\,-1,\,1,\,-1),\quad
\ket{u_3}=\tfrac12(1,\,i,\,-1,\,-i).
\]

Confira \(\ket{u_1}\): \(U_7\) move cada coeficiente uma casa (\(\ket1\to\ket7\) etc.), dando \(\tfrac12(i,\,1,\,-i,\,-1)=i\cdot\tfrac12(1,\,-i,\,-1,\,i)\). Logo \(U_7\ket{u_1}=i\ket{u_1}=e^{2\pi i\cdot\frac14}\ket{u_1}\): fase \(\tfrac14\). Do mesmo jeito, \(\ket{u_s}\) tem fase \(s/4\). No caso geral, com período \(r\):
\[
\ket{u_s}=\frac1{\sqrt r}\sum_{k=0}^{r-1}e^{-2\pi i\,sk/r}\ket{a^k \bmod N},\qquad U_a\ket{u_s}=e^{2\pi i\,s/r}\ket{u_s},\quad s=0,\dots,r-1.
\]
\textbf{A fase de \(U_a\) é \(s/r\), e o período está escondido no denominador.}

\textbf{O truque.} Não sabemos fabricar \(\ket{u_s}\) (isso exigiria conhecer \(r\)), mas somando os quatro:
\[
\ket{u_0}+\ket{u_1}+\ket{u_2}+\ket{u_3}=\tfrac12\,(4,\,0,\,0,\,0)=2\ket{1},\qquad\text{ou seja}\qquad \frac1{\sqrt r}\sum_{s=0}^{r-1}\ket{u_s}=\ket{1},
\]
porque as fases se cancelam em todos os termos, exceto em \(k=0\). E \(\ket1\) é fácil de preparar. Pela seção ``vários autoestados'' da estimação de fase, rodar o algoritmo com \(\ket1\) no alvo sorteia um \(s\) ao acaso (cada um com probabilidade \(1/r\)) e entrega a estimativa de \(s/r\).

Para \(N=15\) e \(a=7\), sai \(s/4\) com \(s\in\{0,1,2,3\}\). Qual é a probabilidade de a medida cair em \(s=1\) ou \(s=3\), que dão \(r=4\) diretamente? (Para \(s=0\) não se aprende nada; para \(s=2\) a fração reduzida é \(\tfrac12\).)

O circuito é a estimação de fase da seção anterior, com \(t=2n\) qubits no registro (basta \(2^t\ge N^2\), para ter precisão suficiente na etapa de frações contínuas; para \(N=21\), \(n=5\), \(t=9\) já serve), controlando \(U_a^{2^k}\) e finalizando com a QFT inversa.

Para \(N=15\) (\(n=4\) bits), com \(t=2n\), quantos qubits tem o registro de fase?

\subsubsection{Das frações contínuas ao período}\label{das-frauxe7uxf5es-contuxednuas-ao-peruxedodo}

Medimos \(y\), com \(y/2^t\approx s/r\). Precisamos recuperar \(r\) a partir de um decimal aproximado. O método é o das \textbf{frações contínuas}: escreve-se o número como \(a_0+\cfrac{1}{a_1+\cfrac1{a_2+\cdots}}\), onde os \(a_i\) são os quocientes do algoritmo de Euclides. Cortar a expansão depois de cada termo dá uma lista de frações cada vez melhores, os \textbf{convergentes}.

Exemplo com \(N=21\), \(a=2\) e \(t=9\) (os restos de \(2^x \bmod 21\) são \(1,2,4,8,16,11,1,\dots\), então \(r=6\)). Uma medida típica é \(y=85\), com \(85/512\approx0{,}1660\). Aplicando Euclides a \(85/512\):

\begin{center}\begin{tabular}{llll}
\hline
passo & divisão & quociente & resto \\
\hline

\hline

1 & \(512 \div 85\) & 6 & 2 \\
2 & \(85 \div 2\) & 42 & 1 \\
3 & \(2 \div 1\) & 2 & 0 \\
\end{tabular}\end{center}

Quocientes: \(0;\,6,\,42,\,2\) (o primeiro, \(0\), é a parte inteira de \(85/512\)). Os convergentes:

\begin{itemize}
\tightlist
\item
  só \(0\): \(\tfrac01\);
\item
  até o \(6\): \(0+\tfrac16=\tfrac16\);
\item
  até o \(42\): \(\cfrac{1}{6+\frac1{42}}=\tfrac{42}{253}\);
\item
  completo: \(\tfrac{85}{512}\).
\end{itemize}

Regra: escolha o \textbf{último convergente com denominador menor que \(N\)}. Aqui é \(\tfrac16\) (o próximo, \(253\), passa de \(21\)), então \(r=6\) é o candidato. Verificamos classicamente: \(2^6=64=3\cdot21+1\equiv1\). Confirmado, e \(r\) é par: \(2^3=8\) dá \(\gcd(7,21)=7\) e \(\gcd(9,21)=3\).

Isso funciona quando \(y\) é o inteiro mais próximo de \(2^t s/r\): então \(|y/2^t-s/r|<\tfrac1{2r^2}\) (pois \(r<N\) e \(2^t\ge N^2\)), e \(s/r\) aparece entre os convergentes. Esse \(y\) sai com probabilidade de pelo menos \(4/\pi^2\) (por \(s\)), então às vezes a medida é ruim ou dá um \(s\) com \(\gcd(s,r)>1\) (a fração reduzida tem denominador menor que \(r\)). Nesses casos o teste \(a^r\equiv1\) falha e \textbf{repete-se}. Em \(N=15\), \(a=7\), o caso \(s=2\) é assim: \(\tfrac24=\tfrac12\), e \(7^2=49\equiv4\ne1\).

Qual é o candidato a \(r\) se \(N=15\), \(t=8\) e a medida foi \(y=192\)?

Implemente a recuperação de \(r\):

\subsubsection{O custo total e o que não é o Shor}\label{o-custo-total-e-o-que-nuxe3o-uxe9-o-shor}

Cada multiplicação modular custa da ordem de \(n^2\) portas e há \(t=2n\) delas, o que dá cerca de \(n^3\) no total (a QFT, com \(O(n^2)\) portas, é mais barata). Só a \textbf{estimação de fase} e a \textbf{QFT} são quânticas; escolha de \(a\), mdc, frações contínuas e teste final ficam no computador clássico. Fatorar números grandes (RSA de 2048 bits) exige milhares de qubits com correção de erros, muito além dos aparelhos atuais.

Qual parte do Shor usa a transformada de Fourier quântica?

\begin{quote}
Para se aprofundar: B. Valiron, \emph{Introduction to Quantum Algorithms and Quantum Programming} (notas de curso, 2025), seção 6.2. Algoritmo original: P. Shor, ``Polynomial-time algorithms for prime factorization and discrete logarithms on a quantum computer'' (1997).
\end{quote}

\subsection{Algoritmo HHL}
Resolver sistemas lineares \(A\vec x=\vec b\) é o cavalo de batalha da ciência computacional, e o tamanho do sistema cresce rápido. O algoritmo HHL (de Harrow, Hassidim e Lloyd) usa a estimação de fase para ``inverter'' a matriz \(A\) \textbf{dentro de um estado quântico}, com custo que cresce só com \(\log N\) no tamanho. Vale entender o que ele promete, e também o que \textbf{não} promete.

\subsubsection{O problema e o que a saída é}\label{o-problema-e-o-que-a-sauxedda-uxe9}

Seja \(A\) uma matriz \(N\times N\) \textbf{hermitiana} (igual à sua adjunta; todos os autovalores são reais) com \(N=2^n\), e \(\vec b\) um vetor. Queremos \(\vec x=A^{-1}\vec b\). Em vez de entregar os números \(x_j\), o HHL entrega um estado de \(n\) qubits com essas amplitudes:

\[
\ket{b}=\frac{1}{\|\vec b\|}\sum_j b_j\ket j,\qquad \ket{x}=\frac{1}{\|\vec x\|}\sum_j x_j\ket j .
\]

Estados quânticos já são normalizados, então perde-se o comprimento de \(\vec x\) (só a direção resta).

\subsubsection{\texorpdfstring{Exemplo \(2\times2\) resolvido à mão}{Exemplo 2\textbackslash times2 resolvido à mão}}\label{exemplo-2times2-resolvido-uxe0-muxe3o}

\[
A=\begin{pmatrix}2&1\\1&2\end{pmatrix},\qquad \vec b=\begin{pmatrix}1\\0\end{pmatrix}\ (\ket b=\ket0).
\]

\(A\) tem autovetores \(\ket+=\tfrac1{\sqrt2}(1,1)\) com autovalor \(\lambda_1=3\), e \(\ket-=\tfrac1{\sqrt2}(1,-1)\) com \(\lambda_2=1\) (confira: \(A\ket+=(3,3)/\sqrt2\)). Escreva \(\ket b\) na base de autovetores:

\[
\ket0=\tfrac1{\sqrt2}\ket+ + \tfrac1{\sqrt2}\ket-.
\]

Inverter \(A\) é dividir cada componente pelo seu autovalor:

\[
A^{-1}\ket0=\tfrac1{\sqrt2}\left(\tfrac13\ket+ + \tfrac11\ket-\right)\ \propto\ \tfrac13\ket++\ket-\ \propto\ (2,-1)^T.
\]

(porque \(\tfrac13(1,1)+(1,-1)=(\tfrac43,-\tfrac23)\).) Normalizando, \(\ket x=\tfrac1{\sqrt5}(2\ket0-\ket1)\). Confira multiplicando \(A\) por \((2,-1)\): dá \((3,0)\), proporcional a \(\vec b\).

A ideia do HHL é fazer exatamente isso com \(N\) gigante: decompor \(\ket b\) em autovetores de \(A\) (sem conhecê-los), dividir cada componente pelo autovalor, recompor.

\subsubsection{Quatro palavras novas}\label{quatro-palavras-novas}

\begin{itemize}
\tightlist
\item
  \textbf{Registro de relógio}: grupo de \(t\) qubits extra, onde a estimação de fase escreve o autovalor (é o ``registro'' da estimação de fase).
\item
  \textbf{Ancila}: um qubit auxiliar, que começa em \(\ket0\) e serve de ``marcador'' do ramo certo.
\item
  \textbf{Descomputar}: aplicar ao contrário uma etapa para apagar o que ela escreveu, e assim desemaranhar o registro de volta a \(\ket0\).
\item
  \textbf{Pós-seleção}: medir a ancila e \textbf{guardar só os resultados} em que deu 1, repetindo o algoritmo quando deu 0.
\end{itemize}

\subsubsection{Os passos, com o exemplo}\label{os-passos-com-o-exemplo}

Para não confundir com o \(t\) da estimação de fase (número de qubits), chamamos de \(\tau\) o \textbf{tempo} em \(e^{iA\tau}\).

\begin{enumerate}
\def\labelenumi{\arabic{enumi}.}
\tightlist
\item
  \textbf{Carregar \(\ket b\).} Prepara-se o registro de dados em \(\ket b=\sum_j\beta_j\ket{u_j}\), onde \(\ket{u_j}\) são os autovetores de \(A\) (desconhecidos) e \(\beta_j=\langle u_j\rangle{b}\). No exemplo, \(\beta_1=\beta_2=\tfrac1{\sqrt2}\).
\item
  \textbf{Estimação de fase de \(e^{iA\tau}\).} A matriz \(e^{iA\tau}\) é unitária (porque \(A\) é hermitiana) e tem os mesmos autovetores de \(A\), com autovalores \(e^{i\lambda_j\tau}=e^{2\pi i\varphi_j}\), onde \(\varphi_j=\lambda_j\tau/2\pi\) é proporcional a \(\lambda_j\). A estimação de fase escreve \(\varphi_j\) (e portanto \(\lambda_j\)) no relógio:
  \[
  \sum_j\beta_j\ket{u_j}\ket{\tilde\lambda_j}.
  \]
  Aqui supomos \(A\) com autovalores positivos e \(\tau\) tal que \(\lambda_{\max}\tau<2\pi\) (assim \(\varphi_j\in(0,1)\); autovalores negativos dariam \(\varphi\) perto de 1 e exigiriam uma leitura com sinal). \textbf{No exemplo}, escolha \(\tau=\pi/2\): então \(\varphi=\lambda/4\), e os autovalores \(3\) e \(1\) viram \(\varphi=\tfrac34\) e \(\tfrac14\), exatos com 2 qubits de relógio (\(\ket{11}\) e \(\ket{01}\)). O estado fica \(\tfrac1{\sqrt2}\big(\ket+\ket{11}+\ket-\ket{01}\big)\).
\item
  \textbf{Rotação controlada pelo relógio.} Para cada valor \(\lambda\) lido no relógio, gira-se a ancila:
  \[
  \ket0\ \longmapsto\ \sqrt{1-\frac{C^2}{\lambda^2}}\,\ket0+\frac{C}{\lambda}\,\ket1 .
  \]
  A constante \(C\) precisa de \(C\le\lambda_{\min}\) para que \(C/\lambda\le1\) seja uma amplitude válida. \textbf{No exemplo}, com \(C=1\): o ramo de \(\ket+\) (\(\lambda=3\)) fica com \(\sqrt{8/9}\ket0+\tfrac13\ket1\), e o de \(\ket-\) (\(\lambda=1\)) fica com \(\ket1\) (sem \(\ket 0\)).
\item
  \textbf{Descomputar.} Aplica-se a estimação de fase \textbf{ao contrário}, o que zera o relógio de volta para \(\ket{00}\) e o desemaranha (apaga-se o ``lixo'' para que ele não atrapalhe a interferência). \textbf{No exemplo}, resta \(\tfrac1{\sqrt2}\big[\ket+(\sqrt{8/9}\ket0+\tfrac13\ket1)+\ket-\ket1\big]\ket{00}\).
\item
  \textbf{Medir a ancila.} Se der \(1\), o registro de dados fica
  \[
  \ket{x}\ \propto\ \sum_j \frac{\beta_j}{\lambda_j}\ket{u_j}=A^{-1}\ket b ,
  \]
  que no exemplo é \(\propto\tfrac13\ket++\ket-\propto(2,-1)\), como na conta à mão. Se der \(0\), tenta-se de novo.
\end{enumerate}

O efeito da rotação foi pôr \(1/\lambda_j\) na amplitude, e a medida da ancila isolou esse ramo.

\subsubsection{Probabilidade de sucesso}\label{probabilidade-de-sucesso}

A probabilidade de a ancila dar 1 é
\[
P_1=\sum_j|\beta_j|^2\frac{C^2}{\lambda_j^2}.
\]
No exemplo com \(C=1\): \(P_1=\tfrac12\cdot\tfrac19+\tfrac12\cdot1=\tfrac59\). Se usássemos \(C=0{,}5\) (ainda \(\le\lambda_{\min}=1\)), qual seria \(P_1\) (4 casas decimais)?

Autovalores muito diferentes pioram isso. O \textbf{número de condição} é \(\kappa=\lambda_{\max}/\lambda_{\min}\) (no exemplo, \(3/1=3\)). Se \(A\) for reescalada para \(\lambda_{\max}=1\), então \(\lambda_{\min}=1/\kappa\) e \(C\approx1/\kappa\), o que pode deixar \(P_1\) da ordem de \(1/\kappa^2\) (reduzida com amplificação de amplitude, a um custo de ordem \(\kappa\) repetições). Quanto vale \(P_1\) para \(A=\mathrm{diag}(1,100)\) (\(\kappa=100\)), \(C=1\) e \(\ket b=\ket1\) (só no autovetor de autovalor 100)?

\subsubsection{Conferindo no computador}\label{conferindo-no-computador}

A parte ``quântica'' do HHL, com \(A\) pequena, pode ser simulada classicamente: decompõe-se \(\ket b\) nos autovetores, divide-se por \(\lambda_j\) e normaliza. Implemente isso. (Use \texttt{np.linalg.eigh}, que devolve autovalores e autovetores de uma matriz hermitiana: \texttt{lamb,\ U\ =\ np.linalg.eigh(A)}, e as colunas de \texttt{U} são os \(\ket{u_j}\). A decomposição de \(\ket b\) é \texttt{beta\ =\ U.T\ @\ b}; a solução é \texttt{U\ @\ (beta\ /\ lamb)}, depois normalizada.)

Agora a probabilidade de sucesso da ancila:

\subsubsection{\texorpdfstring{Limitações: o que o HHL \textbf{não} resolve}{Limitações: o que o HHL não resolve}}\label{limitauxe7uxf5es-o-que-o-hhl-nuxe3o-resolve}

O custo cresce com \(\log N\) nos qubits, mas as letras miúdas pesam:

\begin{itemize}
\tightlist
\item
  \textbf{Leitura.} A saída é o estado \(\ket x\), não a lista \(x_1,\dots,x_N\). Medir só dá uma amostra por execução, e reconstruir os \(N\) números custa pelo menos da ordem de \(N\) execuções, o que anula o ganho. O HHL serve quando queremos uma \textbf{grandeza global} de \(\ket x\) (como \(\bra x M\ket x\), para um observável \(M\), ou comparar com outro estado), não a solução inteira.
\item
  \textbf{Entrada.} Preparar \(\ket b\) também precisa ser barato; carregar um vetor arbitrário de \(N\) números custa da ordem de \(N\) operações, a não ser que \(\vec b\) tenha estrutura (ou venha de um passo quântico anterior).
\item
  \textbf{Esparsidade.} Implementar \(e^{iAt}\) de forma eficiente exige que \(A\) seja \textbf{esparsa} (poucos elementos não nulos por linha, digamos \(s\)) e que se saiba calcular onde estão e quanto valem. O custo cresce com \(s\).
\item
  \textbf{Condicionamento.} O custo cresce com \(\kappa\) (na versão original, com \(\kappa^2\); versões posteriores melhoraram para quase linear em \(\kappa\)), e se \(A\) for quase singular (\(\lambda_{\min}\) perto de 0) o HHL fica caro. Também é preciso um relógio com resolução suficiente para separar \(\lambda_{\min}\) de 0.
\end{itemize}

Quando a matriz é esparsa e bem condicionada e só se quer uma propriedade da solução, o ganho pode ser exponencial. Para matrizes de \textbf{posto baixo} com acesso por amostragem, porém, algoritmos clássicos ``dequantizados'' eliminam a vantagem exponencial; para esparsas e bem condicionadas não se espera isso. Por isso se diz que o HHL é uma \textbf{sub-rotina} promissora, e não uma solução geral para sistemas lineares.

O que torna a leitura de todas as componentes de \(\ket x\) cara no HHL?

\begin{quote}
Para se aprofundar: B. Valiron, \emph{Introduction to Quantum Algorithms and Quantum Programming} (notas de curso, 2025), seção 6.3. Artigo original: A. Harrow, A. Hassidim, S. Lloyd, ``Quantum algorithm for linear systems of equations'' (2009).
\end{quote}

